% ============================================================
\section{Theoretical Background}
\label{sec:theory}
% ============================================================

\subsection{Neutron Footprint Geometry}

The neutron footprint of a CRNS probe is characterised by two
length scales: the radial sensitivity radius $\rzz$ and the
penetration depth $\zzz$.
Both depend on the total hydrogen content in the environment, which
includes liquid soil water $\thetav$, clay-mineral lattice water $\lw$,
and soil organic matter.

\paragraph{Footprint radius.}
\citet{kohli2015} showed that the radius enclosing 86\% of the
neutron signal is well approximated by
\begin{equation}
  \rzz(\thetav, P) = \frac{A_1}{\thetav + A_2} \cdot \frac{P_0}{P},
  \label{eq:r86}
\end{equation}
with $A_1 = \SI{29.13}{\metre\cubic\metre\per\cubic\metre}$,
$A_2 = \SI{0.0578}{\cubic\metre\per\cubic\metre}$, reference pressure
$P_0 = \SI{1013.25}{\hecto\pascal}$, and site pressure $P$.
At $\thetav = \SI{0.20}{\cubic\metre\per\cubic\metre}$ and sea level,
$\rzz \approx \SI{113}{\metre}$.
The footprint shrinks by roughly 40\% from dry ($\thetav \to 0$) to
saturated ($\thetav \approx \SI{0.45}{\cubic\metre\per\cubic\metre}$)
conditions.

\paragraph{Penetration depth.}
The depth enclosing 86\% of the neutron signal is
\begin{equation}
  \zzz(\thetav, \rho_b, \lw) =
    \frac{\SI{8.3}{\centi\metre}}
         {\rho_b \bigl(\num{0.0564} + \thetav + \lw\bigr)},
  \label{eq:z86}
\end{equation}
where $\rho_b$ is the dry soil bulk density [\si{\gram\per\cubic\centi\metre}],
$\thetav$ is liquid volumetric water content [\si{\cubic\metre\per\cubic\metre}],
and $\lw$ is the lattice water content [\si{\gram\per\gram}]
\citep{bogena2013, kohli2021}.
Including $\lw$ in Eq.~\eqref{eq:z86} is physically necessary because
clay-mineral-bound hydrogen moderates neutrons even in air-dry soil,
reducing $\zzz$ regardless of liquid water.

\paragraph{Radial sensitivity weight.}
The contribution of a soil volume element at radial distance $r$ from
the probe decreases exponentially \citep{kohli2015}:
\begin{equation}
  W(r) = \exp\!\left(-\frac{r}{\rzz / 3}\right).
  \label{eq:Wr}
\end{equation}
This weight function is used throughout the pipeline to compute
footprint-averaged quantities (TWI, LULC fractions, soil properties).

% ---------------------------------------------------------------
\subsection{The Desilets Inversion Formula}
\label{sec:desilets}
% ---------------------------------------------------------------

The relationship between neutron count rate $N$ and liquid soil
moisture $\thetav$ is described by the Desilets equation
\citep{desilets2010}:
\begin{equation}
  \frac{N}{N_0} = a_0 + \frac{a_1}{\thetav + \lw + a_2},
  \label{eq:desilets}
\end{equation}
with parameters $a_0 = 0.0808$, $a_1 = 0.372$, $a_2 = 0.115$
\citep{desilets2010}.
The term $\lw$ has been added here (K\"ohli et al., 2021): it represents
the constant offset from lattice water, which is indistinguishable from
liquid water by the neutron probe.
The formula retains full analytical invertibility:
\begin{equation}
  \thetav = \frac{a_1}{N/N_0 - a_0} - \lw - a_2.
  \label{eq:desilets_inv}
\end{equation}
This is the key advantage of Eq.~\eqref{eq:desilets} over more
complex parameterisations \citep{kohli2015}: $N_0$ can be determined
unambiguously from a single calibration measurement.

\paragraph{Reference count $N_0$.}
$N_0$ is computed from the theoretically expected epithermal neutron
count at the site under ``dry'' conditions ($\thetav = 0$, only $\lw$
present):
\begin{equation}
  N_0 = \frac{N_\mathrm{neut,site}}{a_0 + a_1\,/\,(lw + a_2)},
  \label{eq:N0}
\end{equation}
where $N_\mathrm{neut,site}$ is the absolute expected count rate after
altitude, rigidity, and topographic corrections (Section~\ref{sec:fluxes}).
With $\lw = 0$, Eq.~\eqref{eq:N0} reduces to the original Desilets
expression.

\paragraph{Supplementary hydrogen inventories.}
Two additional hydrogen pools can bias $N_0$ estimates but are treated
as supplementary information rather than modifications to
Eq.~\eqref{eq:desilets}:
(i)~soil organic carbon contributes an effective hydrogen equivalent
$\theta_\mathrm{SOC} \approx 0.55\,\rho_b\,\omega_\mathrm{SOM}$
[\si{\cubic\metre\per\cubic\metre}] \citep{bogena2017}; and
(ii)~above-ground biomass water equivalent
$\mathrm{AGBH} \approx \mathrm{LAI} \times
\SI{150}{\gram\per\metre\squared}$
\citep{baatz2015, schron2017}.

% ---------------------------------------------------------------
\subsection{Topographic Correction \texorpdfstring{$\kT$}{kappa\_topo}}
\label{sec:kappa_topo}
% ---------------------------------------------------------------

A CRNS probe on flat, homogeneous terrain sees a reference soil volume
$V_0 = \pi \rzz^2 \zzz$.
Real terrain deviates from this ideal, either by presenting extra soil
(terrain above sensor level) or reduced soil (valleys, cliffs).
The topographic correction is defined as
\begin{equation}
  \kT = \frac{V_\mathrm{eff}}{V_0},
  \label{eq:kappa_topo_def}
\end{equation}
where $V_\mathrm{eff}$ is the $W(r)$-weighted effective soil volume
actually present.

The pipeline implements a \emph{3-D ray-casting} algorithm.
For each ray emitted from the sensor at zenith angle $\theta$ and
azimuth $\phi$, the algorithm computes:
\begin{enumerate}[nosep]
  \item The air path length $L_\mathrm{air}$ from sensor to the first
        DEM intersection, giving the attenuation factor
        $\exp(-L_\mathrm{air}/\lambda_\mathrm{air})$ with
        $\lambda_\mathrm{air} = \SI{130}{\gram\per\centi\metre\squared}$
        / $\rho_\mathrm{air}$.
  \item The soil path length $L_\mathrm{soil} = \zzz/\cos\alpha$,
        corrected for the local surface slope angle $\alpha$ via the
        DEM gradient, giving the production term
        $1-\exp(-L_\mathrm{soil}/\lambda_\mathrm{soil})$ with
        $\lambda_\mathrm{soil} = \SI{162}{\gram\per\centi\metre\squared}$
        / $\rho_b$.
  \item The angular weight $\cos^2\theta$, accounting for both the
        Lambert emission pattern and the solid-angle element
        $\mathrm{d}\Omega = \cos\theta\,\mathrm{d}\theta\,\mathrm{d}\phi$.
\end{enumerate}
The ray contribution is
\begin{equation}
  \mathrm{d}N \propto
    \exp\!\Bigl(-\tfrac{L_\mathrm{air}}{\lambda_\mathrm{air}}\Bigr)
    \Bigl[1-\exp\!\Bigl(-\tfrac{L_\mathrm{soil}}{\lambda_\mathrm{soil}}\Bigr)\Bigr]
    \cos^2\theta\,,
  \label{eq:ray}
\end{equation}
and $\kT$ is the ratio of the sum over all DEM rays to the sum over
an equivalent flat reference surface.
$\kT > 1$ indicates excess terrain (over-estimation of SM without
correction); $\kT < 1$ indicates missing terrain (under-estimation).

% ---------------------------------------------------------------
\subsection{Muon Field-of-View Correction \texorpdfstring{$\kM$}{kappa\_muon}}
\label{sec:kappa_muon}
% ---------------------------------------------------------------

High-energy cosmic muons arriving at the surface contribute to the
secondary neutron production via spallation reactions in soil.
Their differential flux follows $\mathrm{d}N/\mathrm{d}\Omega
\propto \cos^n\theta_z$ with $n \approx 2$.
Surrounding terrain that rises above the horizon at elevation angle
$\psi(\phi)$ blocks all muons with zenith angle
$\theta_z > 90^\circ - \psi(\phi)$, reducing the effective muon flux.

The muon correction is defined as the ratio of the visible-sky
integral to the full-hemisphere integral:
\begin{equation}
  \kM = \frac{\displaystyle\int_0^{2\pi}
              \int_0^{\theta_{z,\max}(\phi)}
              \cos^2\theta_z \sin\theta_z\,\mathrm{d}\theta_z\,\mathrm{d}\phi}
             {\displaystyle\int_0^{2\pi}
              \int_0^{\pi/2}
              \cos^2\theta_z \sin\theta_z\,\mathrm{d}\theta_z\,\mathrm{d}\phi},
  \label{eq:kappa_muon}
\end{equation}
where $\theta_{z,\max}(\phi) = \pi/2 - \psi(\phi)$.
The azimuthal integral evaluates analytically per bin:
\begin{equation}
  f(\phi) = \frac{1 - \cos^3\theta_{z,\max}(\phi)}{3}.
  \label{eq:kappa_muon_az}
\end{equation}
By construction $\kM \leq 1$; it equals 1 on a flat site with no
surrounding topography.

% ---------------------------------------------------------------
\subsection{Land-Use/Land-Cover Correction \texorpdfstring{$\kL$}{kappa\_lulc}}
\label{sec:kappa_lulc}
% ---------------------------------------------------------------

Non-soil hydrogen sources within the footprint — vegetation water,
standing water, built structures — modulate neutron moderation.
Following \citet{schron2017} and \citet{iwema2017}, each land-cover
class $c$ is assigned a hydrogen equivalence factor $f_H^{(c)}$,
normalised to bare soil at the reference moisture $\thetav^\mathrm{ref}$.
The correction is
\begin{equation}
  \kL = \frac{\displaystyle\sum_{i\in\mathrm{footprint}}
              W(r_i)\,f_H^{(c_i)}\,A_i}
             {\displaystyle\sum_{i\in\mathrm{footprint}}
              W(r_i)\,A_i},
  \label{eq:kappa_lulc}
\end{equation}
where $A_i$ is the pixel area and $c_i$ its WorldCover class.
Selected $f_H$ values are listed in Table~\ref{tab:fH}.

\begin{table}[h]
\centering
\caption{Hydrogen equivalence factors $f_H$ by land cover class
(relative to bare soil baseline $f_H = 1$).
Sources: \citet{schron2017}, \citet{iwema2017}, \citet{franz2012}.}
\label{tab:fH}
\begin{tabular}{lcc}
\toprule
Land-cover class      & $f_H$ & Notes \\
\midrule
Water / wetland       & 4.0 / 2.5 & $\approx 4\times$ soil H \\
Snow / glacier        & 3.0       & \\
Tree cover (forest)   & 1.5       & includes canopy water \\
Shrubland             & 1.15      & \\
Grassland / cropland  & 0.95–1.00 & \\
Built-up              & 0.05      & concrete, asphalt \\
Bare / sparse veg.    & 0.03      & \\
\bottomrule
\end{tabular}
\end{table}

For the OSM vector layer the baseline contribution of uncovered
footprint area is computed as
$f_H^\mathrm{base}\times(\Sigma_\mathrm{DEM} - \Sigma_\mathrm{OSM})$
where both sums are $W(r)$-weighted integrals, avoiding the
spatially averaged approximation used in simpler implementations.

% ---------------------------------------------------------------
\subsection{Site Flux Estimates and \texorpdfstring{$N_0$}{N0}}
\label{sec:fluxes}
% ---------------------------------------------------------------

The expected epithermal neutron count rate at the site is
\begin{equation}
  N_\mathrm{neut,site} =
    N_\mathrm{sl}\,f_{R_c}\,
    \exp\!\bigl[\beta_n\,(P_0 - P)\bigr]\,\kT,
  \label{eq:Nneut_site}
\end{equation}
where $N_\mathrm{sl}$ is the detector-specific sea-level reference
count, $f_{R_c} = \exp[\alpha_\mathrm{H}(R_c^\mathrm{site}
- R_c^\mathrm{ref})]$ is the geomagnetic rigidity correction
\citep{hawdon2014} with $\alpha_\mathrm{H} = -0.075\,\si{\per\giga\volt}$,
and $\beta_n = \SI{0.0077}{\per\hecto\pascal}$ is the
barometric coefficient for neutrons.
The site pressure $P$ follows the standard exponential atmosphere
with scale height $H = \SI{8500}{\metre}$.
Rigidity cutoff $R_c$ is retrieved from the Smart \& Shea (2019)
tables via \texttt{crnpy}.

The full correction chain is
\begin{equation}
  \kTot = \kT \cdot \kM \cdot \kL,
  \label{eq:kappa_tot}
\end{equation}
and $N_0$ follows from Eq.~\eqref{eq:N0}.
